\section{A hypergeometric primitive and finite polylogarithmic transport}
\label{dp:section}

\subsection{A coefficientwise regularized \texorpdfstring{${}_4F_3$}{4F3}}
The centered prefactor has a useful consequence for antiderivatives.
For the same parameter domain, set
\begin{equation}\label{dp:Q}
 q_n(u)=\frac{r_n(u)}{n+\eta},\qquad
 Q_u(z)=\sum_{n=0}^{\infty}q_n(u)z^n,\qquad |z|<1.
\end{equation}
Generically this is the Gamma-normalized hypergeometric function
\begin{align}
 Q_u(z)={}&
 \frac{\Gamma(a)\Gamma(b)\Gamma(c)\Gamma(d_0-u)}
 {\Gamma(1+a-b)\Gamma(1+a-c)\Gamma(b+c+u)}\notag\\
 &\times{}_4F_3\!\left(\begin{matrix}
 a,b,c,d_0-u\\1+a-b,1+a-c,b+c+u
 \end{matrix};z\right).
 \label{dp:4F3}
\end{align}
At a nonpositive lower parameter, the product in \eqref{dp:4F3} is
defined by the coefficient series \eqref{dp:Q}, or equivalently by
regularizing the lower Gamma factors before evaluation. The ordinary
${}_4F_3$ and its prefactor must not be separately evaluated as an
undefined expression times zero. The definition \eqref{dp:Q} retains
the zero summands and all their parameter derivatives.

Let $\theta=z\,\mathrm d/\mathrm dz$ and use the Lerch transcendent
\cite[Section 25.14]{dlmf-lerch}
\[
 \Phi(z,s,\eta)=\sum_{n=0}^{\infty}\frac{z^n}{(n+\eta)^s},
 \qquad |z|<1.
\]
The subtracted generating function is
\begin{equation}\label{dp:F}
 F_{K,u}(z)=(\theta+\eta)Q_u(z)
        -\sum_{j=0}^{K-1}C_j(u)\Phi(z,1+2u+2j,\eta).
\end{equation}
It is exactly the generating function of the brackets in
\eqref{dg:RK}. For $-K<\Re u<d_0$, their absolute summability gives
\begin{equation}\label{dp:endpoint}
 \lim_{z\uparrow1}F_{K,u}(z)=R_K(u).
\end{equation}
The limit and every fixed spectral derivative are locally uniform on
compact subsets of this strip. The identity follows from the
summable majorant in Theorem~\ref{dg:master}, not merely from an
unqualified appeal to boundary continuation.

\subsection{An exact weighted primitive}
\begin{theorem}[The ${}_5F_4$-to-${}_4F_3$ primitive]\label{dp:primitive}
For $|z|<1$ on a fixed branch of $z^\eta$, the normalized radial
primitive of \eqref{dp:F} is
\begin{align}
 \int_0^z t^{\eta-1}F_{K,u}(t)\dd t
  =z^\eta\left\{Q_u(z)-\sum_{j=0}^{K-1}
            C_j(u)\Phi(z,2+2u+2j,\eta)\right\}.
 \label{dp:primitiveeq}
\end{align}
It vanishes at zero. Every fixed spectral derivative and every
parameter derivative within the domain may be taken through this
identity, including the dependence of the weight on $\eta=a/2$.
\end{theorem}
\begin{proof}
The defining series converge normally for $|z|<1$ on compact parameter
sets. Integrating the coefficient of $z^n$ uses
\[
 \int_0^z t^{n+\eta-1}\dd t=\frac{z^{n+\eta}}{n+\eta}.
\]
Thus the coefficient $r_n(u)=(n+\eta)q_n(u)$ loses its centered
factor, while the order of every Lerch coefficient increases by one.
This proves \eqref{dp:primitiveeq}. The condition $\eta>0$ makes
the integrand locally integrable at zero and the displayed primitive
zero there. On slightly enlarged compact parameter sets, additional
logarithmic factors from differentiating $t^{\eta-1}$ are still
integrable. Normal convergence and dominated integration therefore
give the derivative assertions.
\end{proof}

More generally, division of coefficients by $(n+\eta)^\ell$ gives
successive applications of the normalized inverse of $\theta+\eta$:
\begin{align}
 F_{K,u}^{\langle\ell\rangle}(z)
 ={}&\sum_{n=0}^{\infty}
       \frac{r_n(u)z^n}{(n+\eta)^\ell}\notag\\
 &-\sum_{j=0}^{K-1}C_j(u)
            \Phi(z,1+2u+2j+\ell,\eta),\qquad \ell\ge0,
 \label{dp:ladder}\\
 (\theta+\eta)F_{K,u}^{\langle\ell+1\rangle}
 ={}&F_{K,u}^{\langle\ell\rangle}.
 \label{dp:ladderrec}
\end{align}
Here $F_{K,u}^{\langle0\rangle}=F_{K,u}$ and the first step is the
lower-order hypergeometric expression in \eqref{dp:primitiveeq}.
All these identities are coefficientwise regular at moving lower
Gamma zeros.

\subsection{Rational shifts give a finite colored-polylogarithm formula}
Let $q\ge1$, $1\le p\le q$, and $\eta=p/q$. Choose a local $q$th
root $w=z^{1/q}$ and put $\omega=e^{2\pi i/q}$. Root-of-unity
filtering gives
\begin{equation}\label{dp:filter}
 \boxed{z^{p/q}\Phi(z,s,p/q)
       =q^{s-1}\sum_{\ell=0}^{q-1}
                     \omega^{-p\ell}\Li_s(\omega^\ell w).}
\end{equation}
To prove it, expand $\Li_s(\omega^\ell w)$ for $|w|<1$ and use
\[
 \frac1q\sum_{\ell=0}^{q-1}\omega^{\ell(k-p)}
 =\begin{cases}1,&k\equiv p\pmod q,\\0,&\text{otherwise}.
   \end{cases}
\]
The surviving positive indices are $k=qn+p$, including $p=q$ with
$n\ge0$, which yields \eqref{dp:filter}. The formula is entire in
the spectral order $s$ for $|z|<1$.

For $m\ge0$, write $\partial_s^m\Li_s$ for order derivatives, not
argument derivatives. Differentiation of \eqref{dp:filter} gives
\begin{align}
 z^{p/q}\partial_s^m\Phi(z,s,p/q)
 ={}&q^{s-1}\sum_{\ell=0}^{q-1}\omega^{-p\ell}
 \sum_{h=0}^m\binom mh(\log q)^{m-h}
              \partial_s^h\Li_s(\omega^\ell w).
 \label{dp:filterjets}
\end{align}
Thus every rational-shift subtraction term and every spectral jet
in \eqref{dp:primitiveeq} is a finite combination of classical
polylogarithms and their order derivatives. At integer resonances the
primitive orders are the even integers $2-2N+2j$; at the half-integer
points they are the odd integers $1-2M+2j$. Nonpositive integer
polylogarithms themselves reduce to rational functions by repeated
application of $z\,\mathrm d/\mathrm dz$ to $z/(1-z)$, while their
order derivatives retain additional information.

For a larger rational shift $\eta=L+p/q$, $L\ge0$ integer,
one must first apply the finite-head identity
\begin{equation}\label{dp:shift}
 z^\eta\Phi(z,s,\eta)
 =z^{p/q}\Phi(z,s,p/q)
       -\sum_{r=0}^{L-1}\frac{z^{r+p/q}}{(r+p/q)^s}.
\end{equation}
This follows by removing the first $L$ terms and is valid for every
$s$. Differentiating it retains the corresponding finite logarithmic
terms. Dropping this finite head would invalidate the claimed
polylogarithmic reduction for $\eta>1$.

For the quartic family $\eta=1/4$, the four colors are $1,i,-1,-i$:
\begin{equation}\label{dp:quarterfilter}
 z^{1/4}\Phi(z,s,1/4)=4^{s-1}
 \{\Li_s(w)-i\Li_s(iw)-\Li_s(-w)+i\Li_s(-iw)\},
 \quad w=z^{1/4}.
\end{equation}
Together with \eqref{dg:quarticshift} and
\eqref{dp:primitiveeq}, this gives a finite Gaussian-colored
polylogarithm description of the counterterms for every resonant
quartic moment and all its harmonic derivatives. It does not furnish
the missing exact relation certificates for the manuscript's
conjectural $S_6$ or $S_8$ evaluations.

\begin{remark}[Boundary values of the primitive]\label{dp:boundary}
In the strip of Theorem~\ref{dg:master}, the bracketed coefficients
in \eqref{dp:F} are absolutely summable. Dividing them by
$n+\eta$ preserves this property. Therefore the primitive itself has
a finite radial endpoint, obtainable from its one combined series.
At a resonant parameter, however, some separate terms on the right of
\eqref{dp:primitiveeq} can be singular at $z=1$. Their combined
expression, or its normally convergent coefficient series, is the
appropriate endpoint object. No separate assignment of values to
divergent polylogarithmic pieces is implicit.
\end{remark}
